% ==========================================================================
% GENERATED by scripts/split_paper.py. Do not edit this file.
%
% Source: imported_papers/grokking/iclr2027_conference_v2.tex
% Part 3: Experimental setup
%
% Edits belong upstream, in the published paper. To change what the
% chapter includes, change imports/grokking.toml and re-run
% the splitter.
% ==========================================================================
\section{Experimental setup}
\label{sec:setup}

\paragraph{Task and network.} We study modular arithmetic, a central task in grokking research (Section~\ref{sec:background}). The operands $a, b \in \mathbb{Z}_p = \{0, 1, \ldots, p-1\}$ are represented by concatenated one-hot vectors, and the target is $c = (a \circ b) \bmod p$ with $\circ \in \{+, -, \times\}$. Our primary task is addition modulo the prime $p = 31$. We use a fully connected ReLU network with two hidden layers of width $256$ and a linear readout; multilayer perceptrons of this kind are standard in grokking studies of modular arithmetic \citep{gromov2023grokking,morwani2024feature}. Unless stated otherwise, circuit measurements and neuron deletion concern the first hidden layer.

\paragraph{Training.} Our reference learning rule is backpropagation with Adam \citep{kingma2015adam}, cross-entropy loss and coupled weight decay $10^{-3}$, which adds the decay term to the gradient before Adam's moment estimates (L2 regularization); Section~\ref{sec:rules} shows that this choice matters. Runs that we compare with each other share the data split and the initialization. We hold out $40\%$ of the operand pairs, train full batch, and apply credit equalization to the first hidden layer; budgets are stated with each experiment (Appendix~\ref{app:repro}). We classify a run as grokking when training accuracy reaches $1.0$ and test accuracy first exceeds $0.95$ at least $500$ epochs later. Results are means over three seeds unless stated otherwise, with standard deviations in parentheses or as error bars.

\subsection{Learning rules}
\label{sec:ruleset}
We distinguish two parts of training. A \emph{credit assignment algorithm}, such as backpropagation, computes a learning signal for each parameter. We call the signal received by a hidden neuron its \emph{credit}. An \emph{optimizer}, such as SGD or Adam, converts these signals into weight updates. We use \emph{learning rule} for the combination of the two and vary each part in turn (Table~\ref{tab:ruleset}).

Credit equalization is the intervention we introduce. It rescales the magnitude of each neuron's incoming-weight gradient before the gradient enters Adam and leaves its direction unchanged (Section~\ref{sec:rules}). The optimizer comparison keeps backpropagation and replaces Adam with LARS, RMSProp, LAMB, Adafactor, Lion, Muon, Shampoo or Nero. DFC replaces the loss gradient by a controller-based signal. During training, a feedback controller drives the network output toward a target, its signal reaches each layer through the transposed Jacobian of the forward network, and each layer learns from the difference between its controlled and its feedforward activity \citep{meulemans2021dfc}. Section~\ref{sec:grokrules} asks which rules grok at all and adds six rules for that purpose. Feedback alignment and direct feedback alignment replace the transposed forward weights of the backward pass by fixed random matrices \citep{lillicrap2016feedback,nokland2016dfa}, and sign-SGD keeps only the sign of each gradient coordinate \citep{bernstein2018signsgd}. Three constructed rules add a component orthogonal to the gradient that is fresh at every step, fixed for the run, or fixed and fading with the training loss (Appendix~\ref{app:generalization}). All six signals are passed to Adam, and the rules that memorize do not enter the circuit comparisons.

\begin{table}[t]
\caption{Learning rules under test, grouped by how they transform the learning signal. Rules that enter only Section~\ref{sec:grokrules} are not listed. Settings are given in Appendix~\ref{app:repro}.}
\label{tab:ruleset}
\begin{center}
\footnotesize
\setlength{\tabcolsep}{4pt}
\begin{tabular}{@{}p{0.34\linewidth}p{0.63\linewidth}@{}}
\toprule
Learning rule & Transformation of the learning signal \\
\midrule
Backpropagation with Adam (reference) & per-parameter adaptive scaling \citep{kingma2015adam} \\
Credit equalization, then Adam & one positive factor per neuron, applied before Adam \\
\midrule
LARS & one factor per layer \citep{you2017lars} \\
LAMB & per-parameter scaling and a layerwise trust ratio \citep{you2020lamb} \\
RMSProp & per-parameter second-moment scaling \citep{tieleman2012rmsprop} \\
Lion & elementwise sign of a momentum signal \citep{chen2023lion} \\
Adafactor & factored second-moment scaling \citep{shazeer2018adafactor} \\
Muon & matrix orthogonalization \citep{jordan2024muon} \\
Shampoo & matrix preconditioning \citep{gupta2018shampoo} \\
Nero & one factor per neuron and weight constraints \citep{liu2021nero} \\
\midrule
DFC & controller-based credit, no gradient rescaling \citep{meulemans2021dfc} \\
\bottomrule
\end{tabular}
\end{center}
\end{table}
